% Nouveau EGA V (cote 25): archive cote 25, site page 21 (PDF page 22).
% typescript. Draft compared: none.
% AI-transcribed by Claude Opus 5.5 (claude-opus-5-5, adaptive thinking, effort high), Message Batches API, batch ega25b (msgbatch_01A7n5YuXubUPNK1zsepi8RG), 2026-11-13; original AI transcription, subsequently checked against both supplied scan parts; see PAGES_CHECKED.tsv and CHECK_DECISIONS.tsv.
\scanpage{25/21}
\mspage{\hand{18}}

La constructibilité se fait par AQT (c'est en fait un\xout{x} oubli au
par.9 , qu'il serait peut-être encore temps de réparer); la partie b)
résulte en principe de 3.3. , sauf que 3.3. a été énoncé de façon trop
particulière, et doit être généralisé en conséquence (la démonstration
donnée étant d'ailleurs essentiellement inchangée). D'ailleurs, il y
a un lapsus dans l'énoncé 4.4. , qui n'est valable tel quel que si f
est quasi-fini; dans le cas général, au lieu de considérer la\xout{ensemble} \inserted[typed]{dimension}
des parties fermées de $X$ resp. des $Y_{\hand{\xi}}$ , il faut regarder la dimension
de leurs images dans P resp. $H_{\hand{\xi}}$ : Redactor demerdetur.

\margin{Passer au No 5.1}
\deleted{\deleted{\hand{4.5.}} \deleted{Dans \illegible{1} \hand{le} présent n$^o$, nous donnons des} conditions \unsure{moyennant} les
\deleted{quelles l'ensemble E envisagé dans 4.1. est ouvert. \xout{Pxxxxxxx} Il}
\deleted{s'agira ici de propriétés P de nature locale sur X resp. $Y_{\hand{\xi}}$ , de sorte}
\deleted{qu'on peut définir l'ensemble U des $y \hand{\in} \underline{Y}$ tels que (si \hand{$\xi$} désigne l'image}
\deleted{de y dans $\hand{\check{}}P$ ) $Y_{\hand{\xi}}$ satisfasse à P en le point y (resp. $G_{\hand{\xi}}$ satisfasse à}
\deleted{P en y). Nous donnerons d'abord des critères pour que U soit ouvert,}
\deleted{en utilisant le par.12. \xout{xxxxxxxxxxxxxxxxxxxx} Comme en tous cas on a}
\[ \deleted{E = P - h(\underline{Y} - U)} , \]
\deleted{i\xout{l}l s'ensuit que si U est ouvert et X propre sur S (donc h propre et}
\deleted{à fortiori fermé} \inserted{car $\underline{Y} \to \unsure{Y}$ propre sur X et $\overset{\hand{\vee}}{P}$ séparé sur S}\deleted{), alors E est également ouvert.}

\deleted{\uline{Proposition} 4.6. Soit F \inserted[typed]{Module} de présentation finie sur X, et supposons}
\deleted{F Cohen-Macaulay relativement à S, i.e. plat relativement à S, et}
\deleted{pour tout $s \hand{\in} S$, $F_s$ sur $X_s$ est Cohen-Macaulay. \xout{Suppxxxxxx} Supposons}
\deleted{de plus que pour tout $s \hand{\in} S$, \xout{ex} toute composante connexe \deleted{Z} de} \hand{$\unsure{S|\zeta} - \unsure{W^{F_s}|W_{F_s}}$}
\inserted{de dim $d$,} \inserted[typed]{\deleted{\xout{exxxitxdimxxZxxx}} et pour}
\deleted{\xout{Spexx} supp $F_s$ , \xout{dexxdimensionxxxxxx}tout $\hand{\xi} \hand{\in} \overset{\hand{\vee}}{P}_s$ , l'image inverse}
\deleted{$Z_{\hand{\xi}}$ de $H_{\hand{\xi}}$ dans $Z_{k(\hand{\xi})}$ \inserted{$(\unsure{S|\zeta|\Sigma} - \unsure{W^{F_s}|W_{F_s}})$} soit de dimension d-1, (i.e. \inserted[typed]{l'image d'}aucune composante}
\deleted{irréductible de \deleted{Z} n'est contenue dans un hyperplan de $\overline{P}_s$\hand{$\check{}$} ). Alors :}
\deleted{a) Le Module \underline{G} sur \underline{Y} est \inserted{Cohen Macaulay, et à fortiori} plat, relativement à $\hand{\check{}}P$ .}
\deleted{\hand{(i) à (iii) de 12.1.1.}}
\deleted{b) Soit P \xout{xxx} l'une des propriétés suivantes pour un module de-}}

\margin{21}

% NOTE: The left bracket directs the paragraph4.5 to No5.1; site22 explicitly calls its cancellation an error.
